\section{Cùng một khối, ba cách cắt}
\label{sec:ch09-cung-mot-khoi}

\index{tiền huấn luyện!khối giống nhau ở ba kiến trúc}%
Chương~\ref{ch:transformer} dựng một encoder-decoder: một chồng khối encoder mà
mỗi khối là self-attention không che cộng mạng truyền thẳng, và một chồng khối
decoder mà mỗi khối là self-attention có che, cộng attention nối sang encoder,
cộng mạng truyền thẳng.

BERT bỏ chồng bên phải. GPT bỏ chồng bên trái. Cái còn lại ở cả hai bên là cùng
một khối, và hình~\ref{fig:ch09-ba-cach-cat} đặt ba thứ cạnh nhau để chỗ ấy nhìn
ra được ngay.

% The picture measures 436.00pt against this book's 441.02pt measure, which is
% 5pt of margin for three columns. Measured with \savebox rather than judged
% from the built page, because a tikzpicture that overflows inside \centering
% raises no Overfull box and the log cannot see it. Recheck this number before
% widening the picture by so much as a node.
\begin{figure}[htbp]
  \centering
  % Ba cach cat cung mot khoi Transformer: encoder-decoder (Vaswani va cong su
% 2017, kien truc chuong 7), encoder-only (BERT) va decoder-only (GPT).
% Trong tam KHONG phai residual/LayerNorm (do la hinh cua chuong 8) ma la
% thanh phan cua khoi - self-attention, co khi them cross-attention, roi FFN -
% va mask nao di kem tung sublayer. Vi vay hinh nay chi ve mot khoi dai dien
% cho moi ngan xep, boc trong khung cham cham "N x", giong quy uoc cua
% hinh chuong 7 va chuong 8.
% Mono-safe va dung lai dung cac style cua hinh chuong 7: diem1 (nen xam nhat)
% cho self-attention khong che, diem2 (net dut) cho self-attention co che,
% diem3 (net cham, nen xam dam hon) cho attention noi sang encoder. So hieu
% 1/2/3 giu nguyen tu hinh chuong 7 de nguoi doc nhan ra day la cung mot hop.
% Hop FFN dung style mac dinh, khong to nen, giong het nhau o ca bon ngan xep,
% va do rong hop (2.8cm) giu nguyen tu hinh chuong 7: do chinh la diem hinh
% nay muon nguoi doc thay ngay, nen khong the thu nho hop de tiet kiem cho.
\begin{tikzpicture}[
  >= Stealth,
  hop/.style     = {draw, rectangle, inner sep = 2pt, font = \scriptsize,
                     align = center},
  so/.style      = {draw, circle, fill = white, minimum size = 3.2mm,
                     inner sep = 0pt, font = \tiny},
  diem1/.style   = {fill = black!10},
  diem2/.style   = {densely dashed, line width = 0.8pt},
  diem3/.style   = {densely dotted, line width = 0.8pt, fill = black!22},
  mui/.style     = {->, line width = 0.6pt, black!75},
  ngucanh/.style = {->, line width = 1.6pt, black},
  khung/.style   = {black!45, densely dotted, line width = 0.5pt},
  nho/.style     = {font = \scriptsize},
  ghichu/.style  = {font = \footnotesize, black!60},
]

  % ================= CỘT 1a: encoder (x = 0) =================
  \draw[mui] (0, -0.35) -- (0, 0);
  \node[hop, diem1, minimum width = 2.8cm, minimum height = 0.7cm]
    (ea1) at (0, 0.35) {Self-attention};
  \node[so] at (ea1.north east) {1};

  \draw[mui] (0, 0.7) -- (0, 1.05);
  \node[hop, minimum width = 2.8cm, minimum height = 0.7cm]
    (ea2) at (0, 1.4) {FFN theo vị trí};

  \draw[mui] (0, 1.75) -- (0, 2.1);

  \begin{scope}[on background layer]
    \draw[khung, rounded corners = 2pt] (-1.5, -0.2) rectangle (1.5, 1.95);
  \end{scope}
  \node[ghichu, anchor = west] at (1.65, 0.875) {N x};

  % ================= CỘT 1b: decoder (x = 3.6) =================
  \draw[mui] (3.6, -0.35) -- (3.6, 0);
  \node[hop, diem2, minimum width = 2.8cm, minimum height = 1.0cm]
    (da1) at (3.6, 0.5) {Self-attention\\có che};
  \node[so] at (da1.north east) {2};

  \draw[mui] (3.6, 1.0) -- (3.6, 1.35);
  \node[hop, diem3, minimum width = 2.8cm, minimum height = 1.0cm]
    (da2) at (3.6, 1.85) {Attention\\encoder-decoder};
  \node[so] at (da2.north east) {3};

  \draw[mui] (3.6, 2.35) -- (3.6, 2.7);
  \node[hop, minimum width = 2.8cm, minimum height = 0.7cm]
    (da3) at (3.6, 3.05) {FFN theo vị trí};

  \draw[mui] (3.6, 3.4) -- (3.6, 3.75);

  \begin{scope}[on background layer]
    \draw[khung, rounded corners = 2pt] (2.1, -0.2) rectangle (5.1, 3.6);
  \end{scope}
  \node[ghichu, anchor = west] at (5.25, 1.7) {N x};

  % mũi tên nặng: đầu ra ngăn xếp encoder đi thẳng vào cross-attention của
  % decoder. Một mũi tên cho cả ngăn xếp là đủ, không vẽ N lần.
  \draw[ngucanh] (0, 2.1) -- (0, 2.25) -- (2.0, 2.25) -- (2.0, 1.85)
    -- (da2.west);

  % ================= CỘT 2: chỉ encoder, BERT (x = 7.45) =================
  \draw[mui] (7.45, -0.35) -- (7.45, 0);
  \node[hop, diem1, minimum width = 2.8cm, minimum height = 0.7cm]
    (ba1) at (7.45, 0.35) {Self-attention};
  \node[so] at (ba1.north east) {1};

  \draw[mui] (7.45, 0.7) -- (7.45, 1.05);
  \node[hop, minimum width = 2.8cm, minimum height = 0.7cm]
    (ba2) at (7.45, 1.4) {FFN theo vị trí};

  \draw[mui] (7.45, 1.75) -- (7.45, 2.1);

  \begin{scope}[on background layer]
    \draw[khung, rounded corners = 2pt] (5.95, -0.2) rectangle (8.95, 1.95);
  \end{scope}
  \node[ghichu, anchor = west] at (9.1, 0.875) {N x};

  % ================= CỘT 3: chỉ decoder, GPT (x = 11.3) =================
  \draw[mui] (11.3, -0.35) -- (11.3, 0);
  \node[hop, diem2, minimum width = 2.8cm, minimum height = 1.0cm]
    (ga1) at (11.3, 0.5) {Self-attention\\có che};
  \node[so] at (ga1.north east) {2};

  \draw[mui] (11.3, 1.0) -- (11.3, 1.35);
  \node[hop, minimum width = 2.8cm, minimum height = 0.7cm]
    (ga2) at (11.3, 1.7) {FFN theo vị trí};

  \draw[mui] (11.3, 2.05) -- (11.3, 2.4);

  \begin{scope}[on background layer]
    \draw[khung, rounded corners = 2pt] (9.8, -0.2) rectangle (12.8, 2.25);
  \end{scope}
  \node[ghichu, anchor = west] at (12.95, 1.025) {N x};

  % ================= nhãn cột =================
  \node[font = \bfseries\scriptsize] at (1.8, -0.9) {Encoder-decoder};
  \node[nho, black!60] at (1.8, -1.35) {chương 7};

  \node[font = \bfseries\scriptsize] at (7.45, -0.9) {Chỉ encoder};
  \node[nho, black!60] at (7.45, -1.35) {BERT};

  \node[font = \bfseries\scriptsize] at (11.3, -0.9) {Chỉ decoder};
  \node[nho, black!60] at (11.3, -1.35) {GPT};

  % ================= chú giải ba điểm attention, số hiệu giữ nguyên từ
  % hình kiến trúc chương 7 =================
  \node[nho, anchor = north, align = left, black!70, text width = 13cm]
    at (5.85, -1.75)
    {1: self-attention không che - nhìn toàn bộ chuỗi trong khối \\
     2: self-attention có che - chỉ được nhìn về sau, cùng kiểu che của
        nhãn 2 ở kiến trúc gốc \\
     3: attention nối sang encoder - chỉ có ở cột encoder-decoder; hộp này
        biến mất hẳn ở hai cột còn lại, đó chính là khác biệt duy nhất giữa
        chúng và ngăn xếp decoder bên trái};

\end{tikzpicture}

  \caption{Cùng một khối, ba cách cắt. Hộp mạng truyền thẳng giống hệt nhau ở cả
  bốn ngăn xếp; self-attention cũng vậy, và cái phân biệt hai phiên bản của nó
  là mask chứ không phải tham số. Cái duy nhất BERT và GPT bỏ đi là attention
  nối giữa hai chồng, vì mỗi bên chỉ còn một chồng.}
  \label{fig:ch09-ba-cach-cat}
\end{figure}

Đọc hình theo chiều ấy thì \enquote{encoder-only} và \enquote{decoder-only} là
tên của hai phép xóa, không phải tên của hai kiến trúc. Và cái phân biệt hai
phép xóa còn lại với nhau nhỏ hơn nữa: nó là mask.

\subsection*{Chỉ khác nhau ở một ma trận tam giác}

\index{tiền huấn luyện!khác biệt encoder-only và decoder-only}%
Mục~\ref{sec:ch07-che} đã dựng mask nhân quả và đo nó bằng phép can thiệp: bỏ
mask đi thì mô hình đạt mất mát huấn luyện \emph{thấp hơn} và điểm bằng không,
vì nó học được cách đọc trộm token nó phải đoán.

BERT giữ nguyên phép đọc trộm ấy. Đó là toàn bộ định nghĩa của
\enquote{bidirectional} trong nhan đề bài: mọi vị trí nhìn thấy mọi vị trí, hai
chiều, không có ma trận tam giác nào chặn lại. GPT giữ mask.

Cùng bộ tham số, cùng phép tính, khác một ma trận nhị phân. Điều đáng nói là hệ
quả: một khi bỏ mask thì mục tiêu huấn luyện của chương~\ref{ch:transformer}
không dùng được nữa, đúng bằng lý do mục~\ref{sec:ch07-che} đo được. Bài BERT
viết câu ấy ra ở mục 3.1:

\begin{quote}
\enquote{Unfortunately, standard conditional language models can only be trained
left-to-right or right-to-left, since bidirectional conditioning would allow
each word to indirectly \enquote{see itself}, and the model could trivially
predict the target word in a multi-layered context.}
\end{quote}

Nên thứ tự nhân quả không phải: chọn kiến trúc rồi chọn mục tiêu. Nó là chọn
mục tiêu rồi kiến trúc theo sau. Mục~\ref{sec:ch09-phat-kien} là chỗ nói về mục
tiêu.

\subsection*{Dựng lại số tham số của các bài}

\index{tiền huấn luyện!dựng lại số tham số}%
Trước khi đi tiếp, một phép kiểm mà tôi nghĩ là rẻ và hóa ra không. Mọi bài
trên đều in số tham số của mô hình nó dựng, và số ấy suy ra được từ chính bảng
siêu tham số cùng bài. Một khối là bốn phép chiếu \(d \times d\), hai ma trận
truyền thẳng \(d \times \dff\) và \(\dff \times d\), hai layer norm. Cộng bảng
embedding, cộng bảng vị trí, nhân số tầng.

\begin{measured}{số tham số của BERT, dựng lại từ chính L, H và A của bài}
\begin{minted}{text}
model       vocab   pooler mlm   rebuilt        paper      ratio
BERT-base   30000   True   False 109,081,344       110M    0.9916
BERT-base   30522   True   False 109,482,240       110M    0.9953
BERT-base   30522   True   True  110,104,890       110M    1.0010
BERT-large  30000   True   False 334,607,360       340M    0.9841
BERT-large  30522   True   False 335,141,888       340M    0.9857
BERT-large  30522   True   True  336,224,058       340M    0.9889
\end{minted}

30000 là con số bài in; 30522 là con số checkpoint phát hành mang.

Đo bằng \code{experiments/ch09\_counts.py} tại tag \code{ch09}.
\end{measured}

BERT-base quay về được: 110M là 109.48M làm tròn, và cột \code{mlm} cho thấy
tính cả đầu ra masked-LM thì nó thành 110.10M, tức là làm tròn theo chiều nào
cũng ra 110M. Không có gì để nói.

\textbf{BERT-large thì không quay về.} 335.14M so với 340M đã in là lệch 1.4\%,
và không cách đọc nào trong sáu cách trên đóng được khoảng ấy. Tôi không biết
340M đếm thêm cái gì, và nói thẳng như vậy chứ không chọn một cách đọc rồi gọi
phần dư là sai số.

Chuyện tương tự ở GPT-2, và ở đây thì số học nói được nhiều hơn.

\begin{measured}{bảng 2 của GPT-2, dựng lại theo chính từ điển mục 2.3 của bài}
\begin{minted}[fontsize=\footnotesize]{text}
model         printed    V=50257 n=1024        gap        gap / d_model
GPT-2 117M       117M   124,439,808    (1.0636)   7,439,808      9687
GPT-2 345M       345M   354,823,168    (1.0285)   9,823,168      9593
GPT-2 762M       762M   774,030,080    (1.0158)  12,030,080      9398
GPT-2 1542M     1542M   1,557,611,200  (1.0101)  15,611,200      9757

model         V that reproduces the printed count
GPT-2 117M        40,570
GPT-2 345M        40,664
GPT-2 762M        40,858
GPT-2 1542M       40,500
\end{minted}

Đo bằng \code{experiments/ch09\_counts.py} tại tag \code{ch09}.
\end{measured}

Cột cuối của bảng trên là chỗ đọc ra. Khoảng cách chia cho \(\dmodel\) gần như
là hằng số, nên nó là một bảng embedding chứ không phải một lỗi về bề rộng hay
độ sâu: bốn hàng thiếu cùng chừng ấy \emph{hàng} embedding. Giải ngược từng
hàng ra kích thước \tn{từ điển}{vocabulary} làm cho con số in đúng, và bốn phép
giải độc lập cho
40,570, 40,664, 40,858 và 40,500, nằm trong 1\% của nhau.

Mục 2.3 của chính bài nói: \enquote{The vocabulary is expanded to 50,257. We
also increase the context size from 512 to 1024
tokens}~\autocite{radford2019gpt2}. Nên bảng 2 in số tham số đếm bằng từ điển
\emph{trước} khi mở rộng, trong khi đoạn văn cách đó vài dòng thông báo việc mở
rộng. Bài không nói vậy; số học nói vậy. Tôi để nguyên chỗ ấy ở dạng nó có:
bốn phép giải, và một câu của bài đặt cạnh.

GPT-1 thì không có gì để dựng lại, vì bài không in số tham số ở đâu cả. Tôi đọc
hết 12 trang để xác nhận điều đó chứ không suy ra từ việc không tìm thấy. Bài
cũng không in kích thước từ điển: nó nói \enquote{a bytepair encoding (BPE)
vocabulary with 40,000 merges}, mà số merge và kích thước từ điển là hai đại
lượng khác nhau. Nếu bạn tra mô hình này ở nơi khác và gặp con số 117M, thì đó
là một phép dựng lại như phép dựng lại ở trên chứ không phải con số của bài: từ
điển 40,478 của checkpoint phát hành cho 116,536,320, và làm tròn ra 117M. Bài
tập cuối chương có một câu hỏi về chỗ đó.

\subsection*{Một bài không đồng ý với chính nó}

\index{tiền huấn luyện!GPT-3 và hai con số tham số}%
GPT-3 là chỗ phép kiểm này trả công.

\begin{measured}{bảng 2.1 của GPT-3, dựng lại cả tám hàng}
\begin{minted}[fontsize=\footnotesize]{text}
model         heads*d_head  d_model  rebuilt             paper     ratio
GPT-3 Small   768           768      125,226,240          0.125B   1.0018
GPT-3 Medium  1024          1024     355,871,744          0.350B   1.0168
GPT-3 Large   1536          1536     760,300,032          0.760B   1.0004
GPT-3 XL      3072          2048     1,315,723,264        1.300B   1.0121
GPT-3 2.7B    2560          2560     2,651,553,280        2.700B   0.9821
GPT-3 6.7B    4096          4096     6,658,404,352        6.700B   0.9938
GPT-3 13B     5120          5140     12,952,938,780      13.000B   0.9964
GPT-3 175B    12288         12288    174,604,259,328    175.000B   0.9977
\end{minted}

Đo bằng \code{experiments/ch09\_counts.py} tại tag \code{ch09}.
\end{measured}

Hàng cuối dựng lại được 174,604,259,328 tham số. Bảng 2.1 của bài in 175.0B.
Bảng D.1 ở phụ lục D của \emph{cùng bài ấy} in 174,600M cho cùng mô hình. Hai
bảng trong một bài không khớp nhau, và phép dựng lại đứng về phía phụ lục, tới
đúng chữ số bảng D.1 làm tròn tới.

Hai cột đầu của bảng cũng đáng nhìn. Sáu trong tám hàng thỏa
\(h \cdot \dk = \dmodel\), tức đúng quan hệ mục 3.2.2 của Vaswani đặt ra và
mục~\ref{sec:ch07-nhieu-dau} đã dùng. Hàng XL in \(\dmodel\) bằng 2048 với 24
head cỡ 128, tích là 3072. Hàng 13B in \(\dmodel\) bằng 5140 với 40 head cỡ
128, tích là 5120. Tôi đọc lại trang 8 ở độ phân giải đầy đủ để loại trừ khả
năng mình đọc nhầm; bài in đúng như vậy. Bảng trên chép nguyên, không sửa.

Không con số nào trong ba chỗ vừa rồi làm thay đổi một luận điểm nào của ba
bài. Tôi để chúng ở đây vì lý do khác: một con số được trích lại đủ nhiều lần
thì không còn ai dựng lại nó. Phép dựng lại ở trên tốn một file Python không
huấn luyện gì, và nó tìm ra ba chỗ.
